\documentclass[10pt]{article}
\usepackage[margin=20mm]{geometry}
\usepackage{amsmath,amssymb,mathtools,fontspec,unicode-math,xurl,hyperref}
\IfFontExistsTF{DejaVu Serif}{\setmainfont{DejaVu Serif}}{\setmainfont{Latin Modern Roman}}
\setmathfont{latinmodern-math.otf}
\setlength{\parindent}{0pt}\setlength{\parskip}{6pt}\sloppy
\title{Initial review of
Fixed and periodic points of the intersection body operator}
\author{}
\date{26 October 2026}
\begin{document}
\maketitle
Abstract: The inspected version is the publisher’s version of record, DOI 10.1007/s00222-025-01342-z. The principal argument was traced through regularity, symmetrization, variation and equality dependencies. One local mathematical inconsistency is established in the auxiliary \(L^p\)-Minkowski discussion: the displayed functional in Section 9.1(4) does not have the asserted stationary points. An exact unit-ball calculation disproves that assertion; a signed replacement restores the Euler–Lagrange equation in the smooth class. No dependency from this defect to Theorem 1.1 was established. This is an initial review with coverage limits, not scientific acceptance.\par
\section*{Identity, version and outcome}
The paper is by Emanuel Milman, Shahar Shabelman and Amir Yehudayoff, Inventiones mathematicae 241 (2025), 509–558 [paper]; the assigned identifier is M25-0081. Locations below use printed journal pagination. PDF pages are one-based in the acquired 50-page file, so printed page 551 is PDF page 43. Actual byte counts and SHA-256 calculations bound this PDF and publisher metadata to the supplied manifest. The journal version was not inferred from an unverified author manuscript.\par
The demonstrated finding is Minor: a peripheral variational assertion is invalid, without a demonstrated dependency to the central classification. The finding is the mathematical failure of stationarity, not a count of typographical slips. No additional substantive defect was established in the checked main proof chain. That conclusion is restricted to the work below and does not certify all results or every external dependency.\par
\section*{Definitions, quantifiers and checked proof chain}
A star body has a continuous positive radial function \(\rho_K\) on \(S^{n-1}\); its intersection body satisfies \(\rho_{IK}(\theta)=|K\cap\theta^\perp|_{n-1}\). Theorem 1.1 concerns all such bodies in \(\mathbb R^n\), \(n\geq3\), satisfying \(I^2K=cK\) for a positive scalar \(c\), and characterizes them as origin-centered ellipsoids. The fixed-point corollary forces a Euclidean ball. Convexity is not an initial hypothesis of the main theorem.\par
All 50 pages of extracted text, including appendix and references, were read. Thirty-one original pages were visually inspected to verify the principal statements, proof transitions, formulas and finding. The coverage ledger records exact pages and read intervals.\par
\subsection*{Regularity and admissible variations}
Appendix A’s regularity iteration was checked for its positive gain: two spherical Radon transforms supply a gain \(n-2>0\). The bounded Sobolev multiplication step and iteration to arbitrary regularity were followed. Positivity uses \(n\geq3\): vanishing of the nonnegative transform at a direction forces vanishing of its integrand on the associated equator; propagation through orthogonal equators then covers the sphere. The same positive-gain argument does not apply when \(n=2\).\par
The full-measure good directions in Theorem 5.7, the multi-graphical definition, volume preservation in Corollary 4.12, uniform Lipschitz control in Proposition 5.8 and radial bounds in Lemma 5.11 were traced into Definition 6.1 and Proposition 6.3. The last radial differentiability step also follows directly from \(r(t)-r_0=O(t)\). Differentiability of its defining function at \((r_0,0)\), with a nonzero radial derivative, gives \[0=a(r(t)-r_0)+bt+o(t),\qquad r'(0+)=-b/a.\] This argument requires the stated pointwise differentiability, not continuously differentiable data on a neighborhood.\par
A possible null-set concern was checked against the cited Lin–Xi framework [linxi]. Finite fibers alone would not transfer a null set of base points to an arbitrary continuous boundary graph. Here the good-direction hypothesis, rectifiable boundary and the area formula in their Theorem 4.1 supply the additional context. The generic graph objection does not establish a defect under the actual hypotheses.\par
\subsection*{First variation and equality}
The self-adjoint Radon calculation in Proposition 6.2 was checked, including the factor \(n-1\), for \[F_c(K)=|IK|-(n-1)c|K|.\] Its sign differs from the auxiliary functional below. Admissibility and preserved volume imply the zero right derivative in Corollary 6.4.\par
The representations \(|IK|=I_0(K)=I_u(K)\), monotonicity of translated rectangle sections and nonnegative difference quotients in Proposition 7.3 were followed. Nonnegativity is essential in the Fatou passage from zero integrated derivative to almost-everywhere equality. Lemma 7.4’s interval calculation was checked for intersection with each pair of opposite facets, retaining strict versus non-strict inequalities. Theorem 7.5 uses a local open cone to exclude the flat section alternative and force a common hyperplane of fiber midpoints. On printed page 545 the original image confirms \((1+\varepsilon)/(1-\varepsilon)<5/4\); an apparent inverse ratio in text extraction is not a finding. The choice \(\varepsilon=1/20\) gives \(21/19<5/4\).\par
The use of connectedness of the punctured \((n-1)\)-dimensional ball in Lemma 7.6 and its restriction \(n\geq3\) were examined. The final Soltan-type characterization in Theorem 7.8 was traced as an external dependency. Soltan’s complete proof was not independently recovered and re-proved; retrieval failures are recorded below. This is a review limit, not evidence against that characterization.\par
\subsection*{Covariance, boundary cases and remaining scope}
For invertible \(T\), the identities \[I(TK)=|\det T|T^{-T}IK,\qquad
 I^2(TK)=|\det T|^{n-2}T(I^2K)\] were checked against \(IB^n=\omega_{n-1}B^n\) and \(I^2B^n=\omega_{n-1}^{\,n}B^n\). Here \(\omega_{n-1}\) is the unit-ball volume in dimension \(n-1\). For an ellipsoid, \(IK=cK\) forces the shape matrix to be scalar.\par
The excluded planar boundary is substantive. The positive nonconstant radial function \(\rho(\theta)=1+\tfrac15\cos(4\theta)\) is invariant under quarter turns. The central line section has length \(2\rho(\theta+\pi/2)\); thus \(IK=2K\) and \(I^2K=4K\) despite the body not being a disk.\par
The other Section 9 applications were read and distinguished from the principal proof. Every cited inequality and equality theorem in them was not independently re-proved. The distributional Fourier reformulation in Corollary 1.6 was not independently formalized in all functional-analytic details. These limits are not findings. No released computational code was identified for this acquired theoretical paper; its absence is not a defect.\par
\section*{Problem 1: failure of the auxiliary stationarity assertion}
\subsection*{Source claim and exact location}
Equation reference: eq:source, eq:bad. On printed page 551, PDF page 43, Section 9.1(4), equation (9.5) is \[h_K^{\,1-p}S_K=c\,d\theta,\qquad p\geq-n,\quad K=-K ,
 \] where \(h_K\) is the support function and \(S_K\) the surface area measure. The immediately following unnumbered display gives \[F_p(K)=\frac1p\int_{S^{n-1}}h_K^p\,d\theta+\frac1c|K|
 \] and asserts that [eq:source] (\#eq:source) is its Euler–Lagrange equation under support-function perturbations. For \(p=0\), the parenthetical replacement is the log integral, again with the positive volume term. Both plus signs were checked on the acquired original page.\par
\subsection*{Complete counterexample}
Take \(n=3\), \(p=2\), \(c=1\) and \(K=B^3\), the unit Euclidean ball. It is compact, convex, origin-symmetric, smooth and strictly convex, with positive support function \(h_K=1\) and area measure \(S_K=d\theta\). Hence \(p\geq-n\) and \(h_K^{1-p}S_K=d\theta=c\,d\theta\), exactly as required.\par
For \(|t|<1/2\), the bodies \(K_t=(1+t)B^3\) form a smooth support-function perturbation, with \(h_{K_t}=1+t>0\), preserving convexity and origin symmetry. The scalar \(c=1\) remains fixed. The claimed stationarity would give zero derivative. Instead, \[F_2(K_t)=2\pi(1+t)^2+\frac{4\pi}{3}(1+t)^3,\qquad
 \left.\frac{d}{dt}F_2(K_t)\right|_{t=0}=8\pi\ne0 .\] This is exact and involves no singularity, limiting convention or numerical approximation. It disproves the stated variational claim even in the smoothest admissible class.\par
\subsection*{General check and precise repair}
For a smooth strictly convex body of positive curvature and positive support function, let \(h_t=h+tq\), with \(q\) smooth and even and \(t\) sufficiently small. The volume variation is \(\delta|K|[q]=\int q\,dS_K\). Thus \[\delta F_p[q]
 =\int_{S^{n-1}}q h^{p-1}\,d\theta+\frac1c\int_{S^{n-1}}q\,dS_K.\] Under [eq:source] (\#eq:source) this is \(2\int_{S^{n-1}}q h^{p-1}\,d\theta\). The admissible parallel-body perturbation \(q=1\) makes it positive. The logarithmic case has the same failure, with \(h^{-1}\) replacing \(h^{p-1}\).\par
Equation reference: eq:repair. The replacement \[\widetilde F_p(K)=-\frac1p\int_{S^{n-1}}h_K^p\,d\theta+\frac1c|K|,
 \qquad
 \widetilde F_0(K)=-\int_{S^{n-1}}\log h_K\,d\theta+\frac1c|K|
 \] instead has \[\delta\widetilde F_p[q]
 =-\int_{S^{n-1}}q h^{p-1}\,d\theta+\frac1c\int_{S^{n-1}}q\,dS_K .\] Stationarity for every even smooth \(q\) is precisely [eq:source] (\#eq:source) in this class. Both measures are even, so restricting to even test functions loses no information. At \(p=-n\), \(\widetilde F_{-n}=|K^\ast|+|K|/c\), consistent with the polar-volume monotonicity orientation invoked there.\par
This repair restores the Euler–Lagrange identity. It does not independently establish the entire nonsmooth \(L^p\)-Minkowski uniqueness application or every equality condition used in it. Author adoption has not been requested.\par
\subsection*{Alternative interpretations and version comparison}
A sign change is a plausible intended repair, but it changes the displayed functional. No volume constraint is stated that would exclude the radial perturbation. Adding a different constrained problem would require its own multiplier argument and would not validate the displayed unconstrained fixed-\(c\) assertion.\par
The arXiv record consulted listed 2408.08171v4, dated 10 June 2025, as its latest version. Its HTML Section 9.1(4) displayed the same positive integral and positive volume terms [arxiv]. This is a focused comparison, not a full equivalence certificate for the journal and arXiv texts. A focused title/erratum search did not identify a correction; this does not prove that none exists.\par
\section*{Dependence, impact and conclusion}
The defect affects the final discussion of recovering a known \(L^p\)-Minkowski application by a general variational method. The central proof uses the different \(F_c\) in Section 6 and is completed in Section 8. No dependency of Theorem 1.1 or its ellipsoid classification on [eq:bad] (\#eq:bad) was established. A Major impact is therefore unsupported by this evidence.\par
Item | Judgment and limit\par
Problem 1 | Minor: exact derivative \(8\pi\) contradicts the claimed stationarity.\par
Repair | Analytically verified Euler–Lagrange identity for [eq:repair] (\#eq:repair) in the specified smooth symmetric class.\par
Core theorem | No additional established defect within the checked chain; external-proof limits remain.\par
Adoption | Not requested; no author response asserted.\par
Scientific decision | Pending; human review records remain empty.\par
\section*{Actual checks, evidence and failures}
The reviewer-authored script `verification/0081/evidence/independent\_checks.py` was executed with Python 3.13.5 and SymPy 1.13.3; exit status zero. Its JSON records the exact derivatives \[\begin{array}{c|cc}
 &\hbox{displayed functional}&\hbox{signed replacement}\\ \hline
 R=1,\ c=1,\ p\ne0&2n\omega_n&0\\
 R=1,\ c=1,\ p=0&2n\omega_n&0\\
 n=3,\ p=2,\ R=1,\ c=1&8\pi&0
 \end{array}\] using \(n\omega_nR^p/p+\omega_nR^n/c\) and its signed counterpart. It also obtained zero for the planar quarter-turn difference and verified \(21/19<5/4\). These symbolic checks corroborate the analytic calculations; they are not a finite-test proof of a classification. No author code was run.\par
Original page images inspected were PDF pages 1–3, 6, 12–16, 20, 24, 27–41, 43, 45–47 and 49. Extraction and rendering commands, outputs and hashes are saved in the evidence directory. The text was read consecutively to its end. Manifest/index parsing was not counted as scientific review.\par
The Lin–Xi author-hosted PDF was consulted for Theorems 3.1 and 4.1 and the good-direction area-formula context. Soltan’s full Section 7 was not recovered: the official download failed certificate verification and additional retrieval timed out. No TLS bypass was used. The failed download is retained in `Soltan\_reference\_retrieval.json` and is not presented as a completed dependency verification.\par
\section*{Supporting checks}
`cc4bf0a5a4115d187b90dde6ed18e8aa55bd161752062d3855fdb31162da72e0`.\par
Its manifest SHA-256 is\par
`48f264aa3b2859f9981888c2a5e5952532b5adfb10d4d33afda7d774102387af`.\par
The authorized portable compiler was actually identified as Tectonic 0.17.0. The report procedure invokes the current self-contained source with `--untrusted` and performs its final pass with `--only-cached`. Companion build metadata records the actual command, logs, compiler hash, full task-local static input inventories before and after compilation, and output PDF hash. Delivery requires matching input inventories and a separate visual inspection of every report page. These facts are recorded in companion metadata rather than presumed from this TeX source.\par
The same review agent performed the mathematical work, independent calculations, source inspection, report creation, compilation and report inspection. Backend equivalence is not independently asserted. No new software was installed, no author code was executed and no human acceptance is implied.\par
9 E. Milman, S. Shabelman and A. Yehudayoff, Fixed and periodic points of the intersection body operator, Invent. Math. 241 (2025), 509–558. DOI: https://doi.org/10.1007/s00222-025-01342-z. Acquired publisher version of record, 50 pages. E. Milman, S. Shabelman and A. Yehudayoff, Fixed and periodic points of the intersection body operator, arXiv:2408.08171v4, 10 June 2025. https://arxiv.org/html/2408.08171v4. Record and Section 9.1(4) consulted 26 October 2026. Y. Lin and D. Xi, Orlicz affine isoperimetric inequalities for star bodies, Adv. Appl. Math. 134 (2022), 102308. Author-hosted manuscript: https://archive.ymsc.tsinghua.edu.cn/pacm\_download/257/11891-Lip-Star.pdf. V. Soltan, Convex solids with hyperplanar midsurfaces for restricted families of chords, Bul. Acad. Stiinte Repub. Mold. Mat. 2(66) (2011), 23–40. External dependency identified; complete proof retrieval unsuccessful.\par
\end{document}
